\documentclass[conference]{IEEEtran}
\IEEEoverridecommandlockouts
\usepackage{cite}
\usepackage{amsmath,amssymb,amsfonts}
\usepackage{algorithmic}
\usepackage{graphicx}
\usepackage{textcomp}
\usepackage{xcolor}
\usepackage{url}
\usepackage[hidelinks]{hyperref}
\def\BibTeX{{\rm B\kern-.05em{\sc i\kern-.025em b}\kern-.08em
    T\kern-.1667em\lower.7ex\hbox{E}\kern-.125emX}}

\begin{document}

\title{Bidirectional Range-Azimuth Upsampling for One-Bit SAR Imaging
\thanks{This work was supported in part by the National Natural Science Foundation of China under Grant 62571337.}
}

\author{\IEEEauthorblockN{1\textsuperscript{st} Given Name Surname}
\IEEEauthorblockA{\textit{dept. name of organization (of Aff.)} \\
\textit{name of organization (of Aff.)}\\
City, Country \\
email address or ORCID}
~\\
\and
\IEEEauthorblockN{2\textsuperscript{nd} Given Name Surname*}
\IEEEauthorblockA{\textit{dept. name of organization (of Aff.)} \\
\textit{name of organization (of Aff.)}\\
City, Country \\
email address or ORCID}
*Corresponding author
~\\
\and
\IEEEauthorblockN{3\textsuperscript{rd} Given Name Surname}
\IEEEauthorblockA{\textit{dept. name of organization (of Aff.)} \\
\textit{name of organization (of Aff.)}\\
City, Country \\
email address or ORCID}
}

\maketitle

\begin{abstract}
One-bit synthetic aperture radar (SAR) imaging greatly reduces echo precision but introduces strong quantization distortion. This paper proposes bidirectional range-azimuth upsampling (BRAU), which distributes a fixed pre-quantization digital upsampling budget over both SAR dimensions instead of concentrating it in range or azimuth alone. Under a common separable random threshold, bidirectional allocations consistently improve reconstruction fidelity over unidirectional allocations across six sampling budgets. Range-compressed spectral analysis further shows reduced off-support and directional leakage for representative bidirectional processing, supporting two-dimensional spectral redundancy as its operating mechanism.
\end{abstract}

\begin{IEEEkeywords}
one-bit SAR imaging, range-azimuth allocation, upsampling, quantization
\end{IEEEkeywords}

\section{Introduction}

Synthetic aperture radar (SAR) forms high-resolution images through coherent processing of echoes collected along a synthetic aperture \cite{cumming2005sar,curlander1991sar}. High-resolution and wide-swath operation produces large raw-data volumes, placing pressure on onboard storage, processing, and downlink resources. Raw-data coding methods such as block-adaptive quantization and predictive coding reduce this burden while retaining sufficient imaging information \cite{kwok1989baq,attema2010fdbaq,magli2003lossy}. One-bit SAR takes this reduction to the extreme by retaining only the signs of the in-phase and quadrature echo components.

Early signum-coded SAR studies established that useful images can be formed from binary echoes and investigated real-time processing, interferometry, phase quantization, and oversampling \cite{franceschetti1991signum,alberti1991time,cappuccino1996processor,fornaro1997interferometry,pascazio1998onebit,franceschetti1999phase}. Related one-bit compressive-sensing research characterized binary acquisition and reconstruction from sign measurements \cite{candes2008intro,boufounos2008onebit,laska2011trust,jacques2013robust}. These ideas were subsequently applied to SAR and radar imaging through model-based and sparsity-aware reconstruction \cite{dong2015map,wang2019enhanced,ge2023sparse}.

More recent one-bit SAR studies have focused on threshold construction and quantization-distortion suppression. Time-varying thresholds recover amplitude information, single-frequency thresholds shift harmful harmonics, and low-precision or learning-based schemes improve reconstructed imagery \cite{demir2018onebit,zhao2019onebit,zhao2021lowprecision,liu2023twobit,si2023dequantization,nie2025fixed}. These methods show that the sampling grid and threshold spectrum jointly influence the distortion entering the SAR imaging chain. However, existing studies generally prescribe one-dimensional upsampling within a particular threshold scheme and do not ask how a fixed digital upsampling budget should be divided between range and azimuth.

Let $R$ and $A$ denote the range and azimuth upsampling factors, respectively, and let $Q_{\mathrm{up}}=R A$ be the total pre-quantization allocation. Under the same $Q_{\mathrm{up}}$, a unidirectional strategy assigns the full budget to one dimension, whereas a bidirectional strategy uses $R>1$ and $A>1$. This paper formulates bidirectional range-azimuth upsampling (BRAU) and evaluates whether dimensional allocation matters for one-bit SAR imaging. The main contributions are as follows:
\begin{itemize}
    \item We formulate fixed-budget range-azimuth upsampling allocation before one-bit quantization, distinguishing bidirectional raw-echo processing from post-reconstruction image interpolation.
    \item We show on 70 paired samples that bidirectional allocations consistently outperform unidirectional allocations under a common separable random threshold, with a near-optimal plateau among several bidirectional factor pairs.
    \item We provide range-compressed spectral evidence showing that bidirectional allocation reduces off-support and directional leakage relative to unidirectional allocation.
\end{itemize}

\section{Method}

\subsection{Fixed-Budget Range-Azimuth Upsampling}

Let $s(\tau,\eta)\in\mathbb{C}$ denote the complex raw SAR echo, where $\tau$ and $\eta$ are range time and azimuth time. Writing $s=I+jQ$, one-bit quantization is applied separately to the in-phase and quadrature components:
\begin{equation}
    y=\operatorname{sign}\!\left(\operatorname{Re}(s+u)\right)
    +j\operatorname{sign}\!\left(\operatorname{Im}(s+u)\right),
\end{equation}
where $u\in\mathbb{C}$ is the quantization threshold. The sign operation removes continuous amplitude information and introduces nonlinear harmonic distortion.

BRAU allocates a fixed digital upsampling budget
\begin{equation}
    Q_{\mathrm{up}}=R\times A
\end{equation}
between range and azimuth. Range-only and azimuth-only allocations are $(Q_{\mathrm{up}},1)$ and $(1,Q_{\mathrm{up}})$, respectively, while a bidirectional allocation satisfies $R>1$ and $A>1$. The $(1,1)$ case is the no-upsampling baseline.

For each allocation, the raw echo is transformed along the selected dimension, centered with \texttt{fftshift}, zero-padded symmetrically, and transformed back to the enlarged grid. The threshold is generated on this grid, followed by one-bit quantization and range compression. When range is enlarged, the range sampling frequency and range-time grid are updated consistently. After range compression, frequency-domain cropping returns the data to the original grid, after which all methods use identical range cell migration correction (RCMC), image formation, region of interest (ROI), and normalization. Following the frequency-domain interpolation procedures adopted in previous one-bit SAR studies \cite{zhao2019onebit,nie2025fixed}, this implementation uses centered FFT zero-padding rather than regenerating the echo through a new SAR simulation. BRAU is therefore a pre-quantization raw-data operation, not image-domain interpolation.

\begin{table}[t]
\caption{BRAU processing pipeline.}
\label{tab:brau_pipeline}
\centering
\scriptsize
\setlength{\tabcolsep}{3pt}
\begin{tabular}{c p{0.72\linewidth}}
\hline
Step & Operation \\
\hline
1 & Select $(R,A)$ under $Q_{\mathrm{up}}=RA$. \\
2 & Apply centered FFT zero-padding along range and azimuth. \\
3 & Generate the threshold on the enlarged echo grid. \\
4 & Quantize the in-phase and quadrature components to one bit. \\
5 & Perform range compression with the matched sampling parameters. \\
6 & Crop the range-compressed data to the original grid in frequency. \\
7 & Apply common RCMC, image formation, ROI, and normalization. \\
\hline
\end{tabular}
\end{table}

\subsection{Threshold Setting}

The main allocation experiment uses a separable random threshold (RT), following the fluctuating-threshold design principle used in one-bit SAR imaging \cite{nie2025fixed}. Its phase is the sum of independent range and azimuth phase vectors:
\begin{equation}
    u_{\mathrm{RT}}(m,n)
    =A_s\hat{\sigma}\exp\!\left[j\bigl(\phi_r(m)+\phi_a(n)\bigr)\right],
\end{equation}
where $\hat{\sigma}=\sqrt{2/\pi}\operatorname{mean}(|s_{\uparrow}|)$ is the scale estimated from the upsampled echo, and $A_s$ controls the threshold amplitude. The fixed-budget results use $A_s=0.6$ as a common comparison setting.

\subsection{Evaluation Protocol}

The experiments use seven SAR datasets representing different land-cover and structural patterns. Ten stratified windows are selected from each dataset, giving 70 paired samples. The ground truth (GT) is reconstructed from the unquantized complex echo using the same range compression, RCMC, image formation, ROI, and normalization as the one-bit results. RT is evaluated with seed 2026 and the common amplitude $A_s=0.6$ for $Q_{\mathrm{up}}\in\{3,4,6,8,9,10\}$. Fractional factors use the nominal budget $Q_{\mathrm{up}}=RA$, with discrete dimensions implemented as $\operatorname{round}(RN_r)$ and $\operatorname{round}(AN_a)$.

Reconstruction fidelity is evaluated using peak signal-to-noise ratio (PSNR) and structural similarity (SSIM) \cite{wang2004ssim}. Equivalent number of looks (ENL) is evaluated on intensity, obtained by squaring the normalized amplitude. For a homogeneous intensity region, ENL is defined as \cite{feng2015tgv}
\begin{equation}
    \mathrm{ENL}=\frac{\mu_I^2}{\sigma_I^2},
\end{equation}
where $\mu_I$ and $\sigma_I^2$ are the mean and variance within a homogeneous ROI. ENL is calculated for 20 samples from the field and port datasets. A $64\times64$ ROI is selected on each GT image and then reused unchanged for every allocation. PSNR and SSIM use all 70 samples.

For each budget, the best bidirectional and unidirectional groups are selected by mean PSNR. Paired Wilcoxon signed-rank tests \cite{wilcoxon1945ranking} are applied to the corresponding 70-sample vectors.

\section{Experiments and Results}

\subsection{Mechanism Analysis Through Spectral Leakage}

One-bit quantization is known to generate nonlinear harmonic distortion \cite{franceschetti1999phase,zhao2019onebit,si2023dequantization}. In the two-dimensional SAR data considered here, this distortion appears along both range and azimuth. The working interpretation of BRAU is that the effective SAR echo occupies a two-dimensional range-azimuth spectral support, whereas a unidirectional allocation enlarges only one spectral axis. In contrast, a bidirectional allocation supplies frequency-domain redundancy in both dimensions and can reduce the overlap between quantization distortion and the effective support. Fig.~\ref{fig:mechanism} compares representative allocations for $Q_{\mathrm{up}}=4$ on the common frequency range $[-2,2]\times[-2,2]$.

\begin{figure}[t]
\centering
\includegraphics[width=\linewidth]{../V4_Experiments_Output/Mechanism/V4_Mechanism_Common4x4.png}
\caption{Range-compressed spectra for representative $Q_{\mathrm{up}}=4$ allocations on a common $4\times4$ normalized-frequency canvas. All panels use the same color scale. R4A1 and R1A4 extend one spectral dimension, whereas R2A2 distributes the extension over both dimensions.}
\label{fig:mechanism}
\end{figure}

For an intermediate matrix $X$, define the spectral energy
\begin{equation}
    S_X(k_r,k_a)=|\mathcal{F}_2\{X\}(k_r,k_a)|^2.
\end{equation}
The support mask $M$ is obtained from a reference spectrum using a relative magnitude threshold $\tau=0.35$. The off-support energy ratio is
\begin{equation}
    \rho_{\mathrm{off}}
    =\frac{\sum S_X(1-M)}{\sum S_X}.
\end{equation}
Directional leakage is computed from the projected spectra
\begin{equation}
    P_r(k_r)=\sum_{k_a}S_X(k_r,k_a),\quad
    P_a(k_a)=\sum_{k_r}S_X(k_r,k_a),
\end{equation}
and the projected masks
\begin{equation}
    M_r(k_r)=\max_{k_a}M(k_r,k_a),\quad
    M_a(k_a)=\max_{k_r}M(k_r,k_a).
\end{equation}
The resulting ratios are
\begin{equation}
    \rho_r=\frac{\sum P_r(1-M_r)}{\sum P_r},\qquad
    \rho_a=\frac{\sum P_a(1-M_a)}{\sum P_a}.
\end{equation}

\begin{table}[t]
\caption{Spectral leakage after range compression for $Q_{\mathrm{up}}=4$. Lower is better.}
\label{tab:mechanism_metrics}
\centering
\scriptsize
\setlength{\tabcolsep}{3.5pt}
\begin{tabular}{c c c c}
\hline
Allocation & Off-support & Range leakage & Azimuth leakage \\
\hline
R1A1 & 0.926008 & 0.471044 & 0.101785 \\
R4A1 & 0.893954 & 0.368013 & 0.070476 \\
R1A4 & 0.901018 & 0.394427 & 0.072251 \\
R2A2 & \textbf{0.892857} & \textbf{0.363750} & \textbf{0.066680} \\
\hline
\end{tabular}
\end{table}

As shown in Table~\ref{tab:mechanism_metrics}, all upsampled groups reduce leakage relative to R1A1, and R2A2 obtains the lowest off-support, range-direction, and azimuth-direction ratios. This representative result supports the interpretation that bidirectional allocation provides two-dimensional spectral redundancy and reduces the overlap between one-bit quantization distortion and the effective SAR support. The metrics are reported at the range-compression stage; no additional improvement is attributed to RCMC because it does not materially change these spectral-energy ratios.

\subsection{Threshold Design}

Prior one-bit SAR studies show that threshold amplitude influences reconstruction quality \cite{demir2018onebit,nie2025fixed}. Fig.~\ref{fig:threshold_design} shows the RT amplitude sensitivity for the best bidirectional allocations at $Q_{\mathrm{up}}=4,6,9$. The SSIM first increases and then decreases with $A_s$. The best tested values shift from $A_s=0.7$ for $Q_{\mathrm{up}}=4$ to $0.8$ for $Q_{\mathrm{up}}=6$ and $1.0$ for $Q_{\mathrm{up}}=9$. The common $A_s=0.6$ used in the main experiment is therefore a controlled setting for allocation comparison rather than a global optimum for every budget.

\begin{figure}[t]
\centering
\includegraphics[width=0.98\linewidth]{../assert/RT_SSIM_bestAs_curve.png}
\caption{Mean SSIM of RT versus $A_s$ for representative bidirectional allocations.}
\label{fig:threshold_design}
\end{figure}

\subsection{Main Fixed-Budget Results}

Fig.~\ref{fig:scene_examples} shows a representative scene using a common normalized-amplitude scale. Both unidirectional allocations improve over one-bit imaging without upsampling, while R2A2 more closely follows the GT in extended boundaries and local scattering texture. The shared $[0,1]$ color limit makes the four panels directly comparable.

\begin{figure}[t]
\centering
\includegraphics[width=\linewidth]{../V4_Experiments_Output/Mechanism/V4_Scene_SharedColorbar.png}
\caption{Representative normalized-amplitude SAR images for $Q_{\mathrm{up}}=4$ with sample-specific PSNR and SSIM relative to GT. GT is self-compared; all panels share the same $[0,1]$ color scale.}
\label{fig:scene_examples}
\end{figure}

\begin{table*}[t]
\caption{Fixed-budget evaluation metrics under separable RT with $A_s=0.6$.}
\label{tab:evaluation_indexes}
\centering
\scriptsize
\renewcommand{\arraystretch}{0.84}
\setlength{\tabcolsep}{8pt}
\begin{tabular}{c c c c}
\hline
Allocation & PSNR & SSIM & ENL \\
\hline
GT & $\infty$ & 1.0000 & 1.0321 \\
R1A1 & 23.0517 & 0.6359 & 1.0094 \\
\hline
\multicolumn{4}{c}{$Q_{\mathrm{up}}=3$} \\
\hline
R3A1 & 25.4031 & 0.7675 & \textbf{1.0203} \\
R1A3 & 25.3326 & 0.7653 & 1.0170 \\
$\mathrm{R}\sqrt{3}\mathrm{A}\sqrt{3}$ &
\textbf{25.6424} & \textbf{0.7811} & 1.0132 \\
\hline
\multicolumn{4}{c}{$Q_{\mathrm{up}}=4$} \\
\hline
R4A1 & 25.7767 & 0.7857 & 1.0202 \\
R1A4 & 25.7258 & 0.7843 & 1.0174 \\
R2A2 & \textbf{26.0623} & \textbf{0.8026} & \textbf{1.0217} \\
\hline
\multicolumn{4}{c}{$Q_{\mathrm{up}}=6$} \\
\hline
R6A1 & 26.1032 & 0.8038 & 1.0161 \\
R1A6 & 26.1415 & 0.8037 & 1.0173 \\
R2A3 & \textbf{26.5585} & 0.8246 & 1.0175 \\
R3A2 & 26.5525 & \textbf{0.8246} & \textbf{1.0204} \\
\hline
\multicolumn{4}{c}{$Q_{\mathrm{up}}=8$} \\
\hline
R8A1 & 26.3618 & 0.8148 & 1.0173 \\
R1A8 & 26.3358 & 0.8136 & 1.0179 \\
R2A4 & 26.8194 & 0.8363 & 1.0183 \\
R4A2 & 26.8066 & 0.8356 & \textbf{1.0187} \\
$\mathrm{R}\sqrt{8}\mathrm{A}\sqrt{8}$ &
\textbf{26.8473} & \textbf{0.8364} & 1.0161 \\
\hline
\multicolumn{4}{c}{$Q_{\mathrm{up}}=9$} \\
\hline
R1A9 & 26.3978 & 0.8171 & 1.0174 \\
R3A3 & \textbf{26.9167} & \textbf{0.8400} & \textbf{1.0181} \\
R9A1 & 26.4482 & 0.8186 & 1.0180 \\
\hline
\multicolumn{4}{c}{$Q_{\mathrm{up}}=10$} \\
\hline
R1A10 & 26.4481 & 0.8198 & \textbf{1.0192} \\
R2A5 & 26.9418 & 0.8423 & 1.0171 \\
R5A2 & \textbf{26.9702} & \textbf{0.8427} & 1.0173 \\
R10A1 & 26.4630 & 0.8206 & 1.0171 \\
\hline
\end{tabular}
\vspace{2pt}

\parbox{0.96\textwidth}{\footnotesize PSNR and SSIM are averaged over 70 paired samples; ENL uses 20 fixed $64\times64$ homogeneous ROIs. Fractional allocations use nominal $Q_{\mathrm{up}}=RA$ and discrete sizes $\operatorname{round}(RN_r)\times\operatorname{round}(AN_a)$.}
\end{table*}

Table~\ref{tab:evaluation_indexes} reports the RT results under the common $A_s=0.6$ setting. R1A1 gives the lowest reconstruction fidelity, whereas a bidirectional group reaches the highest PSNR and SSIM at every listed budget. For $Q_{\mathrm{up}}=3$, $\mathrm{R}_{\sqrt{3}}\mathrm{A}_{\sqrt{3}}$ improves on the best unidirectional allocation by 0.2393~dB and 0.0137 SSIM. At $Q_{\mathrm{up}}=4$, R2A2 obtains 26.0623~dB and 0.8026. At $Q_{\mathrm{up}}=6$, R2A3 and R3A2 are nearly indistinguishable, with PSNR values of 26.5585 and 26.5525~dB.

The same plateau appears at $Q_{\mathrm{up}}=8$. The fractional balanced allocation $\mathrm{R}_{\sqrt{8}}\mathrm{A}_{\sqrt{8}}$ reaches 26.8473~dB and 0.8364 SSIM, only 0.0279~dB above R2A4 and with nearly identical SSIM. R3A3 is best at $Q_{\mathrm{up}}=9$ with 26.9167~dB and 0.8400 SSIM, while R5A2 is best at $Q_{\mathrm{up}}=10$ with 26.9702~dB and 0.8427 SSIM. Across all six budgets, paired comparisons between the mean-best bidirectional and unidirectional groups give PSNR gains of 0.2393--0.5072~dB, with $p_{\mathrm{PSNR}}\leq4.38\times10^{-8}$ and $p_{\mathrm{SSIM}}\leq4.05\times10^{-13}$.

ENL varies less across allocations and is reported as a descriptive speckle statistic rather than a primary ranking criterion. Its maximum coincides with the fidelity-optimal group at $Q_{\mathrm{up}}=4$ and 9, but not at every budget. The consistent PSNR and SSIM results indicate that bidirectionality, rather than exact balance, is the main allocation property under the separable RT.

\section{Conclusion}

This paper studied how a fixed pre-quantization digital upsampling budget should be divided between range and azimuth in one-bit SAR imaging. Under a common separable random threshold, bidirectional allocations consistently achieved higher reconstruction fidelity than concentrating the same budget in one dimension, while neighboring bidirectional factor pairs formed a near-optimal plateau. Range-compressed spectral measurements further showed that bidirectional processing reduced off-support and directional leakage. BRAU therefore provides a fixed-budget allocation principle that exploits the two-dimensional spectral redundancy of SAR echoes before one-bit quantization.

\IEEEtriggeratref{19}
\begin{thebibliography}{00}

\bibitem{cumming2005sar}
I. G. Cumming and F. H. Wong, \textit{Digital Processing of Synthetic Aperture Radar Data: Algorithms and Implementation}. Norwood, MA, USA: Artech House, 2005.

\bibitem{curlander1991sar}
J. C. Curlander and R. N. McDonough, \textit{Synthetic Aperture Radar: Systems and Signal Processing}. New York, NY, USA: Wiley, 1991.

\bibitem{kwok1989baq}
R. Kwok and W. T. K. Johnson, ``Block adaptive quantization of Magellan SAR data,'' \textit{IEEE Trans. Geosci. Remote Sens.}, vol. 27, no. 4, pp. 375--383, Jul. 1989, doi: 10.1109/36.29557.

\bibitem{attema2010fdbaq}
E. Attema, C. Cafforio, M. Gottwald, P. Guccione, A. Monti Guarnieri, F. Rocca, and P. Snoeij, ``Flexible dynamic block adaptive quantization for Sentinel-1 SAR missions,'' \textit{IEEE Geosci. Remote Sens. Lett.}, vol. 7, no. 4, pp. 766--770, Oct. 2010, doi: 10.1109/LGRS.2010.2047242.

\bibitem{magli2003lossy}
E. Magli and G. Olmo, ``Lossy predictive coding of SAR raw data,'' \textit{IEEE Trans. Geosci. Remote Sens.}, vol. 41, no. 5, pp. 977--987, May 2003, doi: 10.1109/TGRS.2003.811556.

\bibitem{franceschetti1991signum}
G. Franceschetti, V. Pascazio, and G. Schirinzi, ``Processing of signum coded SAR signal: Theory and experiments,'' \textit{IEE Proc. F, Radar Signal Process.}, vol. 138, no. 3, pp. 192--198, Jun. 1991, doi: 10.1049/ip-f-2.1991.0025.

\bibitem{alberti1991time}
G. Alberti, G. Franceschetti, G. Schirinzi, and V. Pascazio, ``Time-domain convolution of one-bit coded radar signals,'' \textit{IEE Proc. F, Radar Signal Process.}, vol. 138, no. 5, pp. 438--444, Oct. 1991, doi: 10.1049/ip-f-2.1991.0057.

\bibitem{cappuccino1996processor}
G. Cappuccino, G. Cocorullo, P. Corsonello, and G. Schirinzi, ``Design and demonstration of a real time processor for one-bit coded SAR signals,'' \textit{IEE Proc.-Radar, Sonar Navig.}, vol. 143, no. 4, pp. 261--267, Aug. 1996, doi: 10.1049/ip-rsn:19960406.

\bibitem{fornaro1997interferometry}
G. Fornaro, V. Pascazio, and G. Schirinzi, ``Synthetic aperture radar interferometry using one bit coded raw and reference signals,'' \textit{IEEE Trans. Geosci. Remote Sens.}, vol. 35, no. 5, pp. 1245--1253, Sep. 1997, doi: 10.1109/36.628791.

\bibitem{pascazio1998onebit}
V. Pascazio and G. Schirinzi, ``Synthetic aperture radar imaging by one bit coded signals,'' \textit{Electron. Commun. Eng. J.}, vol. 10, no. 1, pp. 17--28, Feb. 1998.

\bibitem{franceschetti1999phase}
G. Franceschetti, S. Merolla, and M. Tesauro, ``Phase quantized SAR signal processing: Theory and experiments,'' \textit{IEEE Trans. Aerosp. Electron. Syst.}, vol. 35, no. 1, pp. 201--214, Jan. 1999, doi: 10.1109/7.745692.

\bibitem{candes2008intro}
E. J. Cand\`es and M. B. Wakin, ``An introduction to compressive sampling,'' \textit{IEEE Signal Process. Mag.}, vol. 25, no. 2, pp. 21--30, Mar. 2008, doi: 10.1109/MSP.2007.914731.

\bibitem{boufounos2008onebit}
P. T. Boufounos and R. G. Baraniuk, ``1-bit compressive sensing,'' in \textit{Proc. 42nd Annu. Conf. Inf. Sci. Syst.}, Princeton, NJ, USA, Mar. 2008, pp. 16--21, doi: 10.1109/CISS.2008.4558487.

\bibitem{laska2011trust}
J. N. Laska, Z. Wen, W. Yin, and R. G. Baraniuk, ``Trust, but verify: Fast and accurate signal recovery from 1-bit compressive measurements,'' \textit{IEEE Trans. Signal Process.}, vol. 59, no. 11, pp. 5289--5301, Nov. 2011, doi: 10.1109/TSP.2011.2162324.

\bibitem{jacques2013robust}
L. Jacques, J. N. Laska, P. T. Boufounos, and R. G. Baraniuk, ``Robust 1-bit compressive sensing via binary stable embeddings of sparse vectors,'' \textit{IEEE Trans. Inf. Theory}, vol. 59, no. 4, pp. 2082--2102, Apr. 2013, doi: 10.1109/TIT.2012.2234823.

\bibitem{dong2015map}
X. Dong and Y. Zhang, ``A MAP approach for 1-bit compressive sensing in synthetic aperture radar imaging,'' \textit{IEEE Geosci. Remote Sens. Lett.}, vol. 12, no. 6, pp. 1237--1241, Jun. 2015, doi: 10.1109/LGRS.2015.2390623.

\bibitem{wang2019enhanced}
X. Wang, G. Li, Y. Liu, and M. G. Amin, ``Enhanced 1-bit radar imaging by exploiting two-level block sparsity,'' \textit{IEEE Trans. Geosci. Remote Sens.}, vol. 57, no. 2, pp. 1131--1141, Feb. 2019, doi: 10.1109/TGRS.2018.2864795.

\bibitem{ge2023sparse}
S. Ge, D. Feng, S. Song, J. Wang, and X. Huang, ``Sparse logistic regression-based one-bit SAR imaging,'' \textit{IEEE Trans. Geosci. Remote Sens.}, vol. 61, 2023, Art. no. 5217915, doi: 10.1109/TGRS.2023.3322554.

\bibitem{demir2018onebit}
M. Demir and E. Er\c{c}elebi, ``One-bit compressive sensing with time-varying thresholds in synthetic aperture radar imaging,'' \textit{IET Radar, Sonar Navig.}, vol. 12, no. 12, pp. 1517--1526, Dec. 2018, doi: 10.1049/iet-rsn.2018.5044.

\bibitem{zhao2019onebit}
B. Zhao, L. Huang, and W. Bao, ``One-bit SAR imaging based on single-frequency thresholds,'' \textit{IEEE Trans. Geosci. Remote Sens.}, vol. 57, no. 9, pp. 7017--7032, Sep. 2019, doi: 10.1109/TGRS.2019.2910284.

\bibitem{zhao2021lowprecision}
B. Zhao, L. Huang, and B. Jin, ``Strategy for SAR imaging quality improvement with low-precision sampled data,'' \textit{IEEE Trans. Geosci. Remote Sens.}, vol. 59, no. 4, pp. 3150--3160, Apr. 2021, doi: 10.1109/TGRS.2020.3014300.

\bibitem{liu2023twobit}
S. Liu, B. Zhao, L. Huang, B. Li, and W. Bao, ``Lightweight SAR: A two-bit strategy,'' \textit{Remote Sens.}, vol. 15, no. 2, Art. no. 310, Jan. 2023, doi: 10.3390/rs15020310.

\bibitem{si2023dequantization}
C. Si, B. Zhao, L. Huang, and S. Liu, ``A convolutional de-quantization network for harmonics suppression in one-bit SAR imaging,'' \textit{IEEE Trans. Geosci. Remote Sens.}, vol. 61, pp. 1--16, 2023, doi: 10.1109/TGRS.2023.3330530.

\bibitem{nie2025fixed}
G. Nie, B. Zhao, Q. Liu, L. Huang, and G. Liao, ``One-bit synthetic aperture radar imaging based on fixed-threshold with slow-time fluctuations,'' \textit{IEEE Trans. Geosci. Remote Sens.}, vol. 63, 2025, Art. no. 5200715, doi: 10.1109/TGRS.2024.3519757.

\bibitem{wang2004ssim}
Z. Wang, A. C. Bovik, H. R. Sheikh, and E. P. Simoncelli, ``Image quality assessment: From error visibility to structural similarity,'' \textit{IEEE Trans. Image Process.}, vol. 13, no. 4, pp. 600--612, Apr. 2004, doi: 10.1109/TIP.2003.819861.

\bibitem{feng2015tgv}
W. Feng, H. Lei, and H. Qiao, ``Synthetic aperture radar image despeckling via total generalised variation approach,'' \textit{IET Image Process.}, vol. 9, no. 3, pp. 236--248, 2015, doi: 10.1049/iet-ipr.2013.0701.

\bibitem{wilcoxon1945ranking}
F. Wilcoxon, ``Individual comparisons by ranking methods,'' \textit{Biometrics Bull.}, vol. 1, no. 6, pp. 80--83, Dec. 1945, doi: 10.2307/3001968.

\end{thebibliography}

\end{document}
